\documentclass[11pt,a4paper]{article}
\usepackage[utf8]{inputenc}
\usepackage[T1]{fontenc}
\usepackage{geometry}
\geometry{a4paper, margin=0.9in, headheight=20pt, headsep=25pt}
\usepackage{xcolor}
\usepackage{hyperref}
\usepackage{titlesec}
\usepackage{enumitem}
\usepackage{tikz}
\usetikzlibrary{shapes.geometric, arrows.meta, positioning, calc, shadows}
\usepackage[most]{tcolorbox}
\usepackage{booktabs}
\usepackage[french]{babel}
\usepackage{helvet}
\renewcommand{\familydefault}{\sfdefault}

% Couleurs Premium
\definecolor{auroracyl}{HTML}{00F5D4} % Tactical Cyan
\definecolor{aurorablue}{HTML}{1A56DB} % Neural Blue Plus Profond
\definecolor{darkbg}{HTML}{0B1120}
\definecolor{glassbg}{HTML}{F1F5F9}

% Style des boîtes de section
\newtcolorbox{sectionbox}[1]{
    colback=white,
    colframe=aurorablue!20,
    boxrule=0.5pt,
    left=15pt,
    right=15pt,
    top=10pt,
    bottom=10pt,
    sharp corners=southwest,
    arc=4pt,
    title=\sffamily\bfseries\color{aurorablue} #1,
    colbacktitle=white,
    attach boxed title to top left={xshift=10pt, yshift=-3mm},
    boxed title style={boxrule=0pt, colback=white}
}

% Style des sections standard
\titleformat{\section}{\large\bfseries\sffamily\color{aurorablue}\clearpage\vspace{1em}}{}{0em}{}[\titlerule]
\titleformat{\subsection}{\normalsize\bfseries\sffamily\color{aurorablue!80}}{}{0em}{}

% Configuration hyperref
\hypersetup{
    colorlinks=true,
    linkcolor=aurorablue,
    urlcolor=aurorablue,
}

\begin{document}

% Header / Titre Premium
\begin{tcolorbox}[
    colback=darkbg,
    colframe=aurorablue!50,
    arc=0mm,
    boxrule=0pt,
    bottomrule=3pt,
    left=20pt,
    right=20pt,
    top=30pt,
    bottom=30pt
]
    \centering
    \textcolor{white}{\sffamily\Huge\bfseries AURORA // RAPPORT DE PROJET} \\
    \vspace{0.8em}
    \textcolor{auroracyl}{\sffamily\large VERSION 2.6 // HYBRID_MULTI_AGENT_INTELLIGENCE} \\
    \vspace{1.5em}
    \textcolor{white!80}{\sffamily\small Système d'Intelligence Hybride pour l'Excellence Tactique sur Valorant}
\end{tcolorbox}

\vspace{3em}

\noindent \textbf{AURORA} est une plateforme de \textit{Decision Intelligence} de pointe :
\begin{itemize}[label=\textcolor{auroracyl}{\tiny\faIcon{check}}]
    \item Elle transforme des flux vidéo bruts en analyses tactiques de niveau professionnel.
    \item Elle combine la vision par ordinateur multimodale et un écosystème multi-agents sophistiqué.
\end{itemize}

\section{Description Générale}
Le projet AURORA repose sur trois piliers fondamentaux :
\begin{itemize}[label=\textcolor{auroracyl}{$\blacktriangleright$}]
    \item \textbf{Analyse Vidéo Multimodale} : Exploitation de \textbf{Gemini 1.5 Pro} pour l'interprétation sémantique des visuels (VOD, streams).
    \item \textbf{Intelligence Multi-Agents} : Pipeline de 5 agents spécialisés collaborant pour l'audit, le conseil et la prédiction.
    \item \textbf{Visualisation Tactique} : Génération automatisée de heatmaps, tracés de trajectoires et Cyber-HUD (interface futuriste).
\end{itemize}

\section{Catégorisation des Composants Clés}

\begin{sectionbox}{Backend \& Orchestration}
    \begin{itemize}[label=\textbullet, leftmargin=15pt]
        \item \texttt{aurora\_backend.py} : Serveur FastAPI central gérant les requêtes et la coordination des agents.
        \item \texttt{valorant\_tactical\_engine.py} : Base de connaissances sur les cartes, les agents et la logique tactique.
        \item \texttt{gemini\_vision\_client.py} : Interface SDK pour l'analyse visuelle via Gemini 1.5 Pro.
    \end{itemize}
\end{sectionbox}

\vspace{1em}

\begin{sectionbox}{Intelligence Multi-Agents}
    \begin{itemize}[label=\textbullet, leftmargin=15pt]
        \item \texttt{aurora\_agents.py} : Implémentation des agents \textbf{CleanerBot} (données) et \textbf{AnalystBot} (KPIs).
        \item \texttt{aurora\_advanced\_agents.py} : Logique des agents \textbf{CoachBot} (stratégie) et \textbf{OracleBot} (prédictions).
        \item \texttt{aurora\_rag\_system.py} : Système de recherche vectorielle (FAISS) basé sur les stratégies VCT.
    \end{itemize}
\end{sectionbox}

\vspace{1em}

\begin{sectionbox}{Interface \& Automatisation}
    \begin{itemize}[label=\textbullet, leftmargin=15pt]
        \item \texttt{index.html} \& \texttt{frontend.js} : Cyber-HUD exploitant des designs en \textit{glass-morphism}.
        \item \texttt{n8n\_integration.py} : Connecteurs pour l'automatisation externe via Discord et YouTube.
    \end{itemize}
\end{sectionbox}

\section{Workflow Global du Système}

Le cycle de transformation de la donnée suit quatre phases majeures :
\begin{enumerate}[label=\protect\tikz[baseline=(char.base)]{\node[shape=circle,fill=aurorablue,inner sep=1.5pt] (char) {\textcolor{white}{\tiny\bfseries \arabic*}};}]
    \item \textbf{Ingestion Tactique} : Extraction des frames clés depuis un VOD ou un lien YouTube via le \texttt{VideoAnalyzer}.
    \item \textbf{Vision Layer} : Identification spatiale et événementielle par Gemini Vision (Map, Agents, Spike).
    \item \textbf{Intelligence Layer} : Calcul de l'OVR, comparaison RAG avec les stratégies pro et évaluation des risques.
    \item \textbf{Output Layer} : Projection des résultats sur le Cyber-HUD et notification via n8n.
\end{enumerate}

\section{Architecture du Flux de Données}

\begin{center}
\begin{tikzpicture}[
    node distance=1.5cm and 2cm,
    block/.style={rectangle, draw=aurorablue!40, fill=white, drop shadow={opacity=0.1}, rounded corners, minimum width=3.5cm, minimum height=1cm, text centered, font=\sffamily\scriptsize\bfseries},
    io/.style={ellipse, draw=auroracyl!60, fill=auroracyl!5, drop shadow={opacity=0.1}, minimum width=3cm, minimum height=0.9cm, text centered, font=\sffamily\scriptsize\bfseries},
    arrow/.style={-{Stealth[scale=1.2, color=aurorablue]}, thick, rounded corners, color=aurorablue!60}
]
    % Colonne Centrale
    \node [io] (input) {Vidéo / YouTube};
    \node [block, below=of input] (backend) {FastAPI Backend};
    \node [block, below=of backend] (vision) {Gemini Vision};
    \node [block, below=of vision] (agents) {Multi-Agent Pipeline};
    \node [block, below=of agents] (viz) {Visualisation Layout};
    \node [io, below=of viz] (hud) {Cyber-HUD / Report};

    % Éléments Latéraux (Pipeline)
    \node [block, right=of agents] (rag) {RAG System};
    \node [block, left=of agents] (cleaner) {Analyst / Cleaner};

    % Connexions
    \draw [arrow] (input) -- (backend);
    \draw [arrow] (backend) -- (vision);
    \draw [arrow] (vision) -- (agents);
    \draw [arrow] (agents) -- (viz);
    \draw [arrow] (viz) -- (hud);

    % Interaction Agents
    \draw [arrow] (agents.west) -- (cleaner.east);
    \draw [arrow] (agents.east) -- (rag.west);

    % Boîte Intelligence
    \draw[dashed, aurorablue!30, line width=1pt] ($(cleaner.north west)+(-0.4,0.4)$ ) rectangle ($(rag.south east)+(0.4,-0.4)$);
    \node[anchor=north, aurorablue, font=\sffamily\tiny\bfseries, fill=white, inner sep=2pt] at ($(agents.north)+(0,0.6)$) {PROJECTION INTELLIGENTE};

\end{tikzpicture}
\end{center}

\vspace{2em}

\begin{tcolorbox}[
    enhanced,
    colback=white,
    colframe=auroracyl,
    boxrule=0.5pt,
    arc=4pt,
    title=\sffamily\bfseries FOCUS : Moteur RAG,
    attach boxed title to top center={yshift=-3mm},
    boxed title style={colback=auroracyl, boxrule=0pt},
    coltitle=white,
    shadow={2pt}{-2pt}{0pt}{black!10}
]
\small
Le système \textbf{RAG} est le cœur stratégique d'AURORA :
\begin{itemize}
    \item Il compare vos performances en temps réel à une base de données de stratégies VCT professionnelles.
    \item Il fournit des conseils exploitables de haut niveau.
\end{itemize}
\end{tcolorbox}

\end{document}